\chapter{Hands-on Labs}

\section{Purpose}

This chapter turns the manual into a workbook. Each lab has a command to run,
something specific to observe, and a small modification to try. The goal is not
to benchmark the runtime. The goal is to connect visible behavior to the
mechanisms described earlier.

\begin{notabene}[Lab style]
Run one lab at a time. After each command, inspect the example source and the
runtime module it exercises. The examples are intentionally small enough that
changing them should be part of learning them.
\end{notabene}

\section{Lab 1: Timers and cooperative polling}

Run:

\begin{lstlisting}[language=bash,caption={Timer examples}]
cargo run --example executor_sleep
cargo run --example executor_timeout
cargo run --example executor_race
\end{lstlisting}

Observe that the executor can drive futures that suspend and resume without
using a third-party runtime. Then read \filepath{src/executor/future.rs} and
\filepath{src/executor/timer.rs}.

Try this:

\begin{enumerate}
    \item Change the sleep duration in \filepath{executor_sleep.rs}.
    \item Add a second spawned task that sleeps for a different duration.
    \item Print which task completes first.
\end{enumerate}

The lesson is that timers do not run in their own thread. The executor records
deadlines, waits until the next interesting event, and wakes tasks whose
deadlines expire.

\section{Lab 2: Wake coalescing}

Run the unit tests that mention repeated wakes:

\begin{lstlisting}[language=bash,caption={Wake tests}]
cargo test repeated_wakes_share_one_ready_queue_entry -- --nocapture
\end{lstlisting}

Then read the task queueing code in \filepath{src/executor/task.rs} and
\filepath{src/executor/task\_set.rs}.

Try this:

\begin{enumerate}
    \item Add a tiny future that calls its waker multiple times before
          returning \texttt{Pending}.
    \item Spawn it and inspect the executor snapshot's ready queue length.
    \item Confirm that repeated wake calls do not create repeated ready queue
          entries for the same task.
\end{enumerate}

The lesson is that a waker requests another poll; it is not permission to flood
the scheduler with duplicate queue entries.

\section{Lab 3: Readiness-driven TCP}

Run:

\begin{lstlisting}[language=bash,caption={TCP readiness examples}]
cargo run --example async_tcp_pair
cargo run --example async_tcp_server
cargo run --example async_tcp_scoped_server
\end{lstlisting}

Observe the server shape: an accept loop receives a stream, a handler future
owns that stream, and shutdown waits for handlers.

Try this:

\begin{enumerate}
    \item Change a handler to write two chunks instead of one.
    \item Add a short sleep between chunks.
    \item Watch how the client still observes ordered bytes even though the
          server task suspends between writes.
\end{enumerate}

The lesson is that readiness is retry-oriented. A writable notification does
not mean every byte has been written; the helper must keep retrying until the
operation is complete.

\section{Lab 4: Scheduling groups}

Run:

\begin{lstlisting}[language=bash,caption={Scheduling group demos}]
cargo run --example scheduling_group_demo -- 1
cargo run --example async_tcp_scheduling_groups
cargo run --example sharded_scheduling_groups
\end{lstlisting}

Observe the difference between task placement and scheduling accounting. The
TCP scheduling example shows which group live handler tasks belong to. The CPU
demo shows virtual runtime and charged poll time changing over time.

Try this:

\begin{enumerate}
    \item Change the shares in \filepath{scheduling\_group\_demo.rs}.
    \item Increase the runtime from one second to three seconds.
    \item Compare the \texttt{executed}, \texttt{charged}, and
          \texttt{vruntime} columns.
\end{enumerate}

The lesson is that scheduling groups are cooperative accounting classes, not
preemptive priorities. Long-running futures still need to yield.

\section{Lab 5: Observability snapshots}

Run:

\begin{lstlisting}[language=bash,caption={Observability examples}]
cargo run --example sharded_observability
cargo run --example shard_local_worker_observability
\end{lstlisting}

Observe the task rows. Focus on \texttt{status}, \texttt{age\_ms},
\texttt{state\_ms}, \texttt{polls}, and \texttt{wait}.

Try this:

\begin{enumerate}
    \item Increase the sleep duration inside the worker tasks.
    \item Add one more worker per shard.
    \item Observe how task count, timer count, and state duration change.
\end{enumerate}

The lesson is that snapshots are owned samples. They are not tracing events.
To build a console-like view, sample repeatedly and compute differences.

\section{Lab 6: Cross-shard submission}

Run:

\begin{lstlisting}[language=bash,caption={Cross-shard examples}]
cargo run --example sharded_submit
cargo run --example sharded_map_reduce
cargo run --example sharded_broadcast
\end{lstlisting}

Observe that each example explicitly names where work should run. Nothing
migrates implicitly.

Try this:

\begin{enumerate}
    \item Submit work from shard 0 to shard 1.
    \item Have the remote future return \texttt{current\_executor\_shard()}.
    \item Assert that the returned shard id is the target shard, not the
          caller's shard.
\end{enumerate}

The lesson is that the remote future is polled remotely, and the result crosses
back as an owned value.

\section{Lab 7: Shard-local state}

Run:

\begin{lstlisting}[language=bash,caption={Shard-local examples}]
cargo run --example shard_local
cargo run --example shard_local_current
cargo run --example shard_local_stoppable_workers
\end{lstlisting}

Observe that local values are accessed on owning shards and that workers can be
stopped cooperatively.

Try this:

\begin{enumerate}
    \item Add a second field to the shard-local value.
    \item Update it inside \texttt{with\_current}.
    \item Return an owned summary instead of a reference.
\end{enumerate}

The lesson is that shard-local access is synchronous. If you need to await,
first produce an owned value and await outside the local-state closure.

\section{Lab 8: CPU placement}

Run:

\begin{lstlisting}[language=bash,caption={CPU placement example}]
cargo run --example sharded_cpu_placement
\end{lstlisting}

On Linux, also try the Docker helper:

\begin{lstlisting}[language=bash,caption={Linux CPU placement through Docker}]
tools/linux-docker.sh cargo run --example sharded_cpu_placement
\end{lstlisting}

Observe whether placement is applied, unsupported, or rejected. The exact
result depends on platform and container CPU permissions.

Try this:

\begin{enumerate}
    \item Request more shards than the container allows CPUs.
    \item Compare best-effort placement with required placement.
    \item Inspect the shard snapshots printed by the example.
\end{enumerate}

The lesson is that CPU placement is an explicit request to the operating
system. It is not a load balancer.

\section{Lab 9: Linux \texttt{io\_uring}}

This lab requires a Linux host or container where \texttt{io\_uring} setup is
allowed.

\begin{lstlisting}[language=bash,caption={io\_uring labs}]
SITAS_DOCKER_IO_URING=1 tools/linux-docker.sh cargo run --example os_uring
SITAS_DOCKER_IO_URING=1 tools/linux-docker.sh cargo run --example os_uring_batch
SITAS_DOCKER_IO_URING=1 tools/linux-docker.sh cargo run --example os_uring_abandon
\end{lstlisting}

Observe the dispatcher snapshots. Focus on pending submissions, tracked
operations, buffered completions, registered wakers, and abandoned operations.

Try this:

\begin{enumerate}
    \item Increase the number of submitted operations in
          \filepath{os\_uring\_batch.rs}.
    \item Drop a future before completion in a small variant of the abandon
          example.
    \item Watch how abandonment and deferred buffer counters change.
\end{enumerate}

The lesson is that completion-based I/O has a different cancellation problem
than readiness-based I/O. The kernel may still own an operation after the
future is gone.

\section{Lab 10: Sharded index build}

Run:

\begin{lstlisting}[language=bash,caption={Sharded index build}]
cargo run --example sharded_index_build -- --records 512 --shards 4 --seed 0xbeef --cleanup
\end{lstlisting}

On Linux with \texttt{io\_uring}:

\begin{lstlisting}[language=bash,caption={io\_uring index build}]
SITAS_DOCKER_IO_URING=1 tools/linux-docker.sh cargo run --example sharded_index_build_uring -- --records 512 --shards 4 --seed 0xbeef --cleanup
\end{lstlisting}

Observe how partition work is spread across shards, then merged into one final
sorted index.

Try this:

\begin{enumerate}
    \item Change the record count.
    \item Change the shard count.
    \item Compare per-shard partition summaries and final verification output.
\end{enumerate}

The lesson is that the runtime's value appears when a problem can be split into
owned shard-local work, then combined through explicit owned results.

\section{After the labs}

After running these labs, return to the walkthrough chapter and read the
modules again. The source is easier to understand after seeing the runtime's
observable behavior. If a lab result is surprising, use the debugging checklist
from the walkthrough chapter before changing code.
